% =============================================================================
% exercises/ejercicios_cap04b.tex
% Ejercicios — Capítulo 4, Sección B
% Temas: Grupos funcionales, nomenclatura y propiedades
% =============================================================================

\section*{Ejercicios --- Secci\'on B: Grupos funcionales}
\addcontentsline{toc}{section}{Ejercicios: Grupos funcionales}
\label{sec:ejercicios_cap04b}

\begin{bloque_ejercicios}[Ejercicios B --- Grupos funcionales y nomenclatura]

\begin{ejercicios}

% ----- B1: Identificación de grupos funcionales ----------------------------
\item Relacione cada compuesto con su grupo funcional y su clase
de compuesto org\'anico:

{%
\centering
\begin{tabular}{llll}
\toprule
\textbf{F\'ormula} & \textbf{Grupo funcional} & \textbf{Clase} &
\textbf{Nombre} \\
\midrule
\ce{CH3-OH}                   & & & \\
\ce{CH3-O-CH3}                & & & \\
\ce{CH3-CHO}                  & & & \\
\ce{CH3-CO-CH3}               & & & \\
\ce{CH3-COOH}                 & & & \\
\ce{CH3-COO-CH2CH3}           & & & \\
\ce{CH3-NH2}                  & & & \\
\ce{H2N-CH2-COOH}             & & & \\
\bottomrule
\end{tabular}
\par
}%

\textit{[Respuesta: metanol/alcohol; \'eter dimet\'ilico/\'eter;
etanal/aldeh\'ido; propanona/cetona; \'acido etanoico/\'ac.
carbox\'ilico; etanoato de etilo/\'ester; metilamina/amina;
glicina/amino\'acido]}

% ----- B2: Nomenclatura de alcoholes ---------------------------------------
\item Con las reglas UIQPA para alcoholes:

\begin{enumerate}[label=\alph*)]
    \item Nombre los siguientes alcoholes:
          \begin{enumerate}[label=\roman*)]
              \item \ce{CH3-CH2-CH2-OH}
              \item \ce{CH3-CH(OH)-CH2-CH3}
              \item \ce{CH3-CH2-CH(OH)-CH2-CH3}
              \item \ce{(CH3)2C(OH)-CH2-CH3}
          \end{enumerate}
    \item Clasif\'iquelos (incisos i--iv) como primarios,
          secundarios o terciarios, seg\'un el carbono que porta el
          grupo \ce{-OH}.
    \item Escriba la f\'ormula estructural condensada del
          2-metil-1-propanol y del 2-butanol.
    \item Explique por qu\'e el etanol hierve a \SI{78}{\celsius}
          mientras el \'eter dimet\'ilico (misma f\'ormula molecular
          \ce{C2H6O}) hierve a \SI{-24}{\celsius}.
\end{enumerate}

\textit{[Respuesta: (a-i) 1-propanol;
(a-ii) 2-butanol;
(a-iii) 3-pentanol;
(a-iv) 2-metil-2-butanol (alcohol terciario)]}

% ----- B3: Éteres -----------------------------------------------------------
\item Usando las reglas de nomenclatura de los \'eteres
(\textbf{Cuadro~\ref{eters}}):

\begin{enumerate}[label=\alph*)]
    \item Nombre los siguientes \'eteres seg\'un el sistema UIQPA
          (nombre alcoxi + alcano):
          \begin{enumerate}[label=\roman*)]
              \item \ce{CH3-O-CH2CH3}
              \item \ce{CH3CH2-O-CH2CH3}
              \item \ce{CH3CH2CH2-O-CH2CH2CH2CH2CH3}
          \end{enumerate}
    \item Escriba la f\'ormula del etoxipropano y del
          metoxibenceno (\ce{C6H5-O-CH3}).
    \item El etanol (\ce{CH3CH2OH}) y el \'eter dimet\'ilico
          (\ce{CH3-O-CH3}) son is\'omeros. ?`Cu\'al de los dos
          es soluble en agua y cu\'al tiene mayor punto de
          ebullici\'on? Explique en t\'erminos de los puentes
          de hidr\'ogeno.
\end{enumerate}

\textit{[Respuesta: (a-i) metoxietano; (a-ii) etoxietano;
(a-iii) propoxipentano;
(c) el etanol es m\'as soluble y hierve m\'as alto por
los puentes de hidr\'ogeno intermoleculares del \ce{-OH}]}

% ----- B4: Aldehídos y cetonas ---------------------------------------------
\item Para los aldeh\'idos y cetonas:

\begin{enumerate}[label=\alph*)]
    \item Nombre los siguientes compuestos:
          \begin{enumerate}[label=\roman*)]
              \item \ce{CH3-CH2-CHO}
              \item \ce{CH3-CH2-CH2-CHO}
              \item \ce{CH3-CO-CH3}
              \item \ce{CH3-CH2-CO-CH2-CH3}
          \end{enumerate}
    \item Escriba la f\'ormula condensada del:
          \begin{enumerate}[label=\roman*)]
              \item 4-metilhexanal
              \item 2-pentanona
              \item 4-metil-3-hexanona
          \end{enumerate}
    \item El formaldeh\'ido (HCHO) tiene importantes aplicaciones
          industriales y biol\'ogicas. Mencione dos usos descritos
          en el cap\'itulo y explique por qu\'e su soluci\'on
          acuosa al 40\% se llama formalina.
    \item ?`Cu\'al es la diferencia estructural entre un aldeh\'ido
          y una cetona? Identifique el grupo carbonilo en ambos.
\end{enumerate}

\textit{[Respuesta: (a-i) propanal; (a-ii) butanal;
(a-iii) propanona (acetona); (a-iv) 3-pentanona]}

% ----- B5: Ácidos carboxílicos ----------------------------------------------
\item Con base en la nomenclatura de \'acidos carbox\'ilicos
(\textbf{Cuadro~\ref{grasos}} y \textbf{Cuadro~\ref{cooh}}):

\begin{enumerate}[label=\alph*)]
    \item Nombre los siguientes \'acidos:
          \begin{enumerate}[label=\roman*)]
              \item \ce{CH3-CH2-CH2-COOH}
              \item \ce{CH3-(CH2)_{14}-COOH}
              \item \ce{CH3-CHCl-COOH}
          \end{enumerate}
    \item Escriba la f\'ormula del \'acido butanoico,
          el \'acido hexanoico y el \'acido acetilsalic\'ilico
          (aspirina).
    \item El \'acido propi\'onico (PM 74) hierve a \SI{141}{\celsius}
          mientras que el $n$-pentano (PM 74) hierve a \SI{35}{\celsius}.
          Explique esta gran diferencia en t\'erminos de las
          interacciones intermoleculares.
    \item ?`Por qu\'e los \'acidos grasos con m\'as de 10 \'atomos
          de carbono son s\'olidos a temperatura ambiente?
          Relacione con su volatilidad y peso molecular.
\end{enumerate}

\textit{[Respuesta: (a-i) \'acido butanoico (but\'irico);
(a-ii) \'acido hexadecanoico (palm\'itico);
(a-iii) \'acido $\alpha$-cloropropanoico]}

% ----- B6: Ésteres ----------------------------------------------------------
\item Los \'esteres tienen importantes aplicaciones como aromas
(\textbf{Cuadro~\ref{ester}}):

\begin{enumerate}[label=\alph*)]
    \item Identifique el \'acido y el alcohol que dan origen
          a cada \'ester:
          \begin{enumerate}[label=\roman*)]
              \item Etanoato de isopentilo (sabor a pl\'atano)
              \item Butanoato de etilo (aroma a pi\~na)
              \item Metanoato de isobutilo (sabor a frambuesa)
          \end{enumerate}
    \item Escriba la reacci\'on general de esterificaci\'on entre un \'acido carbox\'ilico \ce{RCOOH} y un alcohol \ce{R}'\ce{OH},  indicando el subproducto formado.
    \item ?`En qu\'e disolventes son solubles los \'esteres? ?`Son
          solubles en agua? Explique con base en su polaridad.
    \item Escriba la f\'ormula del etanoato de octilo y su
          aroma asociado seg\'un el cap\'itulo.
\end{enumerate}

% ----- B7: Aminas -----------------------------------------------------------
\item Las aminas son derivados del amon\'iaco:

\begin{enumerate}[label=\alph*)]
    \item Clasifique las siguientes aminas como primarias,
          secundarias o terciarias e ind\'iquelo:
          \begin{enumerate}[label=\roman*)]
              \item \ce{CH3-NH2}
              \item \ce{(CH3)2NH}
              \item \ce{(CH3)3N}
              \item \ce{C6H5-NH2} (anilina)
          \end{enumerate}
    \item Con los datos del \textbf{Cuadro~\ref{aminas}}, compare
          los puntos de ebullici\'on de metilamina, dimetilamina y
          trimetilamina. ?`Existe una tendencia clara? Explique en
          t\'erminos de puentes de hidr\'ogeno y simetr\'ia
          molecular.
    \item Nombre los siguientes compuestos:
          \begin{enumerate}[label=\roman*)]
              \item \ce{CH3-CH2-NH2}
              \item \ce{(CH3-CH2)_2-NH}
              \item \ce{(CH3-CH2)_3-N}
          \end{enumerate}
    \item Las aminas de bajo peso molecular tienen olores
          desagradables que recuerdan al pescado. ?`Por qu\'e
          las aminas de alto peso molecular ($>$C$_6$) tienen un
          olor menos intenso? Relacione con la volatilidad.
\end{enumerate}

\textit{[Respuesta: (a) \ce{CH3NH2}: primaria;
\ce{(CH3)2NH}: secundaria;
\ce{(CH3)3N}: terciaria;
anilina: primaria arom\'atica;
(c-i) etilamina; (c-ii) dietilamina; (c-iii) trietilamina]}

% ----- B8: Aminoácidos y proteínas -----------------------------------------
\item Los $\alpha$-amino\'acidos son la base de las prote\'inas:

\begin{enumerate}[label=\alph*)]
    \item Escriba la estructura general de un $\alpha$-amino\'acido
          indicando: el carbono alfa, el grupo amino, el grupo
          carboxilo y el grupo variable R.
    \item Para los siguientes amino\'acidos, identifique el grupo R:
          glicina (R = H), alanina (R = \ce{CH3}) y
          metionina (R = \ce{CH3SCH2CH2-}).
    \item ?`Qu\'e es un enlace pept\'idico? Escriba la reacci\'on
          de formaci\'on del diglicina (Gli-Gli) a partir de dos
          mol\'eculas de glicina, indicando el subproducto
          eliminado.
    \item De los 20 amino\'acidos que usa la vida para sintetizar
          prote\'inas, ?`cu\'antos se consideran esenciales para
          el ser humano y por qu\'e?
\end{enumerate}

% ----- B9: Nomenclatura integradora ----------------------------------------
\item Nombre cada uno de los siguientes compuestos indicando
su clase y grupo funcional:

\begin{enumerate}[label=\alph*)]
    \item \ce{CH3-CH2-CH2-CH2-COOH}
    \item \ce{CH3-CH2-CONH2}
    \item \ce{CH3-CH2-CH2-OH}
    \item \ce{CH3-CH2-CO-CH2-CH3}
    \item \ce{CH3-CH2-NH-CH3}
    \item \ce{HCOOCH2CH3}
\end{enumerate}

\textit{[Respuesta:
(a) \'acido pentanoico / \'ac. carbox\'ilico;
(b) propanamida / amida;
(c) 1-propanol / alcohol primario;
(d) 3-pentanona / cetona;
(e) N-metiletilamina / amina secundaria;
(f) metanoato de etilo / \'ester]}

% ----- B10: Escritura de fórmulas ------------------------------------------
\item Escriba la f\'ormula estructural condensada de los
siguientes compuestos:

\begin{enumerate}[label=\alph*)]
    \item \'acido butanoico
    \item 2-metil-3-pentanona
    \item etoxibutano (etil butil \'eter)
    \item 2-etilhexanal
    \item butil propilamina (N-propilbutilamina)
    \item 4-metiloctanamida
\end{enumerate}

\textit{[Respuesta:
(a) \ce{CH3CH2CH2COOH};
(b) \ce{CH3CH(CH3)COCH2CH3};
(c) \ce{CH3CH2-O-CH2CH2CH2CH3};
(d) \ce{CH3CH2CH(CH2CH3)(CH2)2CHO};
(e) \ce{CH3CH2CH2CH2-NH-CH2CH2CH3};
(f) \ce{CH3CH2CH(CH3)CH2CH2CH2CH2CONH2}]}

\end{ejercicios}

\end{bloque_ejercicios}
% =============================================================================
% FIN DE ejercicios_cap04b.tex
% =============================================================================
